\setcounter{section}{23}
\section{Empleo Público}

Nombre completo: Empleo público en España. Tipología del empleo público. Estatuto Básico del Empleado Público. Planificación de recursos humanos en las Administraciones Públicas: la oferta de empleo público. Ley de incopmatibilidades del personal al servicio de las Administraciones Públicas.

% Creado en 2026-06
\localtableofcontents

\subsection{Empleo público en España: Marco general} % 3
    El empleo público constituye un activo para que la Administración Pública cumpla con el mandato del Art. 103.1 de la CE, de servir a los intereses generales. El personal es heterogéneo tanto por la diversidad de administraciones como por los regímenes jurídicos que regulan su relación.

    Existen diferentes típos de empleados según el vínculo jurídico: funcionarios de carrera e interinos, personal laboral y personal eventual. También se organizan en diferentes ámbitos como la administración civil, FFAA, FCSE, personal docente o personal estatutario.

    El empleo público mantiene un volumen elevado y estable, superando los 3,1 millones de efectivos, principalmente en comunidades autónomas y entidades locales con un 82\%.

    Los principales retos se encuentran en la transformación digital, las nuevas tecnologías, la automatización de procesos y la ciberseguridad. Otros retos son el envejecimiento de las plantillas, necesitando atraer talento, la modernización de la selección y mejora de la gestión pública.

    \subsubsection{Modelo español de empleo público} % 3
        El modelo se configura como un sistema mixto, coexistiendo personal funcionario con personal laboral. Cumpliendo la necesidad de equilibrar la independencia con la flexibilidad operativa. Se asienta sobre los pilares del Art.103 de CE que establece las bases del estatuto de los funcionarios públicos. El marco legal vigente es el Real Decreto Legislativo 5/2015, el Estatuto Básico del Empleado Público. Este permite la existencia de personal laboral para funciones no reservadas a funcionarios públicos.

        Estos regímenes están sujetos a un sistema de incompatibilidades y un régimen disciplinario propio, orientado a evitar conflictos de intereses y a asegurar la objetividad de las actuaciones. Esto está regulado en dos regímenes jurídicos:
        \begin{description}
            \item[Derecho Administrativo (Personal Funcionario)] Su rasgo mas relevante es la inamovilidad, no como un privilegio individual, sino como garantía de la neutralidad política. Blinda a los funcionarios  públicos las potestades que implican participación en el ejercicio de las potestades públicas o salvaguardia de los intereses generales.
            \item[Derecho Laboral (Personal Laboral)] Regulado por la legislación laboral común y convenios colectivos, aplicable a puestos técnicos o instrumental
        \end{description}

        El sistema se organiza de forma jerarquizada y profesionalizada:
        \begin{description}
            \item[Sistemas selectivos] Por principios de igualdad mérito, capacidad y publicidad (Arts. 14 y 23.2 de CE y Art. 55 de EBEP) mediante concurso oposición (y excepcionalmente concurso)
            \item[Ordenación] Agrupándose en Cuerpos y Escalas según titulación y funciones, en lugar de puestos de trabajo concretos. Instrumentándose ambos a través de RElaciones de Puestos de Trabajo (Art 74 EBEP).
            \item[Descentralización] Presentando un alto grado de descentralización territorial, siendo su mayoría en CCAA, y Entidades Locales, debido a su competencia directa en la gestión de servicios públicos esenciales (Sanidad, Educación, etc)
        \end{description}

        El modelo se está modernizando bajo el marco de la \textbf{Agenda España Digital 2026}  y el \textbf{Plan de Digitalización de las Administraciones Públicas} hacia la denominada \textbf{Administración 4.0}. Enfocado desde el RDL 6/2023 introduciendo herramientas de gestión avanzada como la evaluación del desempeño, la planificación estratégica para el relevo generacional y la incorporación de tecnologías disruptivas.

    \subsubsection{Evolución histórica de la función pública en España} % 4
        La trayectoria refleja la evolución de la Administración decimonónica basada en el favor político hacia un modeloburocrático profesionalizado, y hacia un modelo moderno de gestión estratégica. Existiendo tres etapas clave:
        \begin{description}
            \item[Fin de las cesantías y profesionalización S.XIX a 1918] Las Cesantías son un mecanismo de rotación del personal según el cambio de gobierno de turno, habiendo una rotación masiva.
            \begin{description}
                \item[Estatutos de Lopez Ballesteros (1825) y Bravo Murillo (1852)] primeras normas con criterios de mérito, categorías profesionales y examenes de contratación.
                \item[Estatuto de Maura 1918] Fin de las cesantías con la inamovilidad de los funcionarios, generalización de la oposición y sistema de escalafón.
            \end{description}
            \item[Modelo Burocrático y la Transición]
            \begin{description}
                \item[Ley de funcionarios civiles del Estado (1964)] Unificó la dispersión normativa previa y consolidó la organización en Cuerpos y Escalas.
                \item[Ley 30/1984 para la reforma de la función pública] Primera gran reforma constitucional, con la nueva estructura autonómica y el fin de los Cuerpos tradicionales. Potenció la figura del personal laboral
            \end{description}
            \item[EBEP a reforma del talento]
            \begin{description}
                \item[Estatuto Básico del Empleado Público] Aprobado en 2007  y unificado en 2015, con conceptos como la gestión de recursos humanos, la carrera horizontal, el personal firectivo profesional y la evaluación del desempeño.
                \item[Ley 20/2021 de medidas urgentes para la reducción de la temporalidad] Para reducir la alta tasa de temporalidad y eimpulso de estabilización con los estándares de la UE
            \end{description}
        \end{description}

        El modelo actual se encuentra en una fase de despliegue profunda con la ley 6/2023. Materializando reformas del \textbf{Plan de Recuperación, Transformación y Resiliencia} con tres transformaciones históricas:
        \begin{description}
            \item[Evaluación del Desempeño y Carrera profesional] Evolucionando de la antigüedad hacia el mérito contínuo. Implantando una evaluación vinculada a la progresión y la percepción de complementos retributivos.
            \item[Planificación Estratégica y Relevo] Abandonando la gestión reactiva en favor de una planificación, priorizando el relevo generacional y la captación de perfiles tecnológicos y especializados.
            \item[Modernización de la selección y Digitalización] Bajo el marco de la \textbf{Agenda España Digital 2026}, transformando los procesos selectivos hacia la agilidad, con la medición de competencias prácticas sobre la memorización.
        \end{description}

    \subsubsection{Based jurídica actual} % 5
        El régimen se articula mediante un bloque jerarquizado que garantiza la profesionalidad, imparcialidad y objetividad. Tras las reformas postpandemia y actualizaciones recientes, el marco legal queda de la siguiente forma.

        A nivel constitucional, la Constitución es la norma de la que emana todo el sistema, asentándose en los siguientes preceptos:
        \begin{description}
            \item[Art. 14] Igualdad de los españoles ante la ley, prohibiendo discriminación en los procesos de selección y promoción
            \item[Art. 23.2] el derecho a acceder en condiciones de igualdad a las funciones y cargos públicos.
            \item[Art. 103.3] Reserva de la ley para el Estatuto de los Funcionarios públicos, el acceso de acuerdo al mérito y capacidad, el derecho de sindiación, las incompatibilidades y la imparcialidad del ejercicio de sus funciones.
            \item[Art. 149.1.18°] La competencia del estado sobre las bases del régimen estatutario de los funcionarios como mínimo común normativo para todos los empleados públicos.
        \end{description}

        A nivel normativo básico y de aplicación general. Normas dictadas al amparo de Art. 149.1.18>
        \begin{description}
            \item[RDL 5/2015 Estatuto Básico del Empleado Público] Define las clases de personal, derechos, deberes y bases comunes.
            \item[L 53/1984 de incpomatibilidades del Personal al Servicio de las Administraciones Públicas] Deducada a garantizar la objetividad, prohibiendo el desempeño de actividades que interfieren con el interés general.
            \item[L 20/2021 de medidas urgentes para la reducción de la temporalidad en el empleo público]  Diseña procesos de estabilización y limita la tasa de interinidad estructural al 8\%
            \item[L 20/1978 de Reconocimiento de Servicios previos en la Administración Pública] Regula el cómputo de la antiguedad en trienios en el sector público
            \item[L 15/2022 para la igualdad de trato y la no discriminación] Refuerza la protección transversal en el acceso y desarrollo de la carrera profesional.
        \end{description}

        Normativa básica que unclute el Estatuto Básico del Empleado Público, la Ley de incompatibilidades del Personal al Servicio de las Administraciones Públicas, la Ley de medidas urgentes para la reducción de la temporalidad del empleo público, ley de reconocimiento de servicios previos de la Administración pública, y Ley integral para la igualdad de trato y la no discriminación.

        Normativa Específica como el plan de Recuperación, Transformación y Resiliencia, la Ley de Medidas para la Reforma de la Función Pública, Decreto de la Ley de Funcionarios Civiles y del Estado y los Reglamentos de Ingreso del Perosnal al servicio de la AGE, de Situaciones Administrtivas de los Funcionarios Civiles
        de Régimen Disciplinario.

    \subsubsection{Datos del empleo público en España} % 7
        Con datos de julio de 2025. El número de efectivos asciende a 3,1 millones. ampliándose con respecto a enero. Siendo el mayor aumento en la Administración local. Actualmente se divide en 540.000 del Estado, casi 2 millones de las comunidades autónomas y 630.000 de la administración local.

        La EPA incluye a empresas y entes intrumentales, y alcanzando 3,6 millones. Siendo una tendencia alcista generalizada.

        Comparado internacionalmente, España se encuentra a menos de un punto por debajo de la OCDE, a 5 de Francia y a 12 del bloque nórdico, pero superando a Alemania en 6 puntos.

        Se está produciendo un rejuvenecimiento del funcionariado, aunque este es ligero. La mayoría del empleo lo realizan mujeres (58,5\%) siendo los sectores mas femeninos las comunidades autónomas, Justicia, Sanidad, Docencia y gestión procesal, y los mas masculinos las Fuerzas Armadas y CSE con un 85\% de hombres.

    \subsubsection{Retos actuales del empleo público} % 13
        Según el informe \textbf{Administración 2030} se sintetizan los retos de la próxima década con las siguientes variables de expresión:
        \begin{description}
            \item[El océano de las incertidumbres] Riesgos sitemáticos, volátiles e impredecibles que exigen una gestión pública adaptativa
            \item[La atomización de la cosa pública] Multiplicación y fragmentación de los actores sociales
            \item[El envejecimiento demográfico] una ciudadanía que presiona los sistemas de protección social
            \item[Exigencia prestacional] Un perfil de ciudadano cualificado que reorienta sus demandas de protección y certidumbre hacia el Estado
            \item[Vectores Socioeconómicos y Territoriales] Polarización en la distribución de la riqueza y empleo.
            \item[Desconfianza institucional y Tsunámi digital] Necesidad de revertir el descrédito crónico con la tecnología como motor de la trasnformación.
        \end{description}

        En cuanto a la integración de la inteligencia artificial y la Gobernanza de datos se ha desarrollado la \textbf{Estrategia Nacional de Inteligencia Artificial} (ENIA) y el \textbf{Reglamento Europeo de la IA} (AI Act). El reto radica en pasar de una burocrácia de trámites a una gobernanza algorítmica. Importante que las decisiones automatizadas respeten los principios de transparencia, no discriminación y supervisión humana (Human In the Loop).

        En cuanto al envegecimiento. El 48\% de los efectivos tienen mas de 55 años, contra el 27\% en la OCD. Con una gran brecha generacional. Ante esta previsible jubilación masiva han aumentado las ofertas de empleo, con una oferta de mas de 37.000 plazas, especialmente en perfiles tecnológicos y en el Cuerpo Superior de Sistemas y Tecnologías de la Información.

        Por los modelos de selección, el proceso memorístico resulta incompleto. Se transiciona hacia la evaluación de competencias, habilidades blandas y la capacidad de aprendizaje conínuo. La ordenación de recursos humanos debe adaptarse para permitir la movilidad funcional hacia sectores críticos con necesidades, o áreas de urgencia social. Esto se plasma en el \textbf{Plan Estratégicco de Recursos humanos de la AGE}

        Relativo a la Ciberseguridad se encuentra la \textbf{Estrategia de Ciberseguridad de la Unión Europea} y las \textbf{Directrices del CCN-CERT}. Siendo el reto el entorno de cibercrimen sofisticado, con el factor humano como eslabón crítico. Siendo necesario implantar una cultura transversal de ciberseguridad, capacitando para la defensa frente a brechas de datos, desinformación y fraude digital. Llevando también la confianza de los ciudadanos en los sistemas públicos.

        El empleo público debe ser parte de una transición verde. Tanto en la reducción propia, como en los criterios de contratación pública y en las políticas diseñadas.


\subsection{Marco normativo del empleo público} % 15
    \subsubsection{Marco constitucional del empleo público} % 16
        La constitución establece unas bases dogmáticas e institucionales para la Administración pública. Se busca que esta sea profesional e imparcial, no como un derecho del funcionariado, sino como una garantía para el correcto funcionamiento del estado.

        Se conecta un principio organizativo con los derechos fundamentales de los ciudadanos y los trabajadores, incluidos los públicos.

        En el Art. 103.3 se indica que \textit{La ley regulará el estatuto de los funcionarios públicos, el acceso a la función pública de acuerdo con los principios de mérito y capacidad, las peculiaridades del ejercicio de su derecho a sindicación, el sistema de incompatibilidades y las garantías para la imparcialidad en el ejercicio de sus funciones}. Derivándose de aquí cuatro mandatos. Una \textbf{Ley Estatutaria} que permita que se mantenga jurídicamente y no por el gobierno de turno. Los principios de \textbf{Mérito y Capacidad} para todo el sistema de seleccin, es decir, la aptitud t€cnica, conocimientos, trayectoria; y la capacidad como idoneidad para el desempeño del puesto. Y la imparcialidad, como forma de blindar al funcionariado ante las presiones políticas o alternancia partidista de los gobiernos.

        En jurisprudencia del Tribunal Constitucional, se han redactado dos sentencias: STC 50/1986 y STC 148/1987 donde se fijan los siguientes parámetros:
        \begin{description}
            \item[Naturaleza del derecho] El Art. 23.2 dice que corresponde al legislador concretar los requisitos, pero otorga a los aspirantes el amparo constitucional ante vulneraciones.
            \item[Derecho de acceso y permanencia] El derecho no se agota al superar el proceso selectivo, ampara el derecho a la permanencia en la función pública y a la progresión, proscribiendo remociones arbitrarias, ceses injustificados o postergaciones.
            \item[Prohibición de requisitos \textit{Ad Personam}] Las convocatorias deben ser generales y abstractas. Se pueden establecer requisitos específicos, pero deben ser proporcionales, razonables y directamente relacionados con el puesto. Siendo lo contrario una reserva de plaza no justificada.
        \end{description}

        El régimen constitucional modula los derechos colectivos de sus efectivos para la defensa de sus intereses laborales y la continuidad de los servicios esenciales de la comunidad.
        \begin{description}
            \item[Libertad Sindical] Se le reconoce a los empleados públicos, pero se permite establecer límites en el ejercicio. Siendo mas estricta en las Fuerzas Armadas y los Institutos Armados de caracter militar (Guardia Civil), a los que se les prohibe el derecho de sindicación, y sre restringe sustancialmente en Cuerpos Policiales.
            \item[Derecho a Huelga] Prácticamente reconocido con jurisprudéncia STC 11/1981, aunque condicionado por la obligación de garantizar el mantenimiento de los servicios esenciales de la comunidad. El incumplimiento puede acarrear sanciones disciplinarias de gravedad.
        \end{description}

    \subsubsection{Estatuto básico del empleado público} % 17
        El Texto Refundido de la Ley del Estatuto Básico del Empleado Público RD 5/2015 (TREBEP) es la piedra angular del sistema de función pública. Como texto refundido, supone arreglar la fragmentación normativa existente. Unificando las bases del régimen estatutario de los funcionarios y las normas específicas aplicables al personal laboral al servicio de las Administraciones públicas.

        El TREBEP se basa en dos artículos de la CE. La competencia del Estado sobre legislación laboral y la competencia para dictar las bases del régimen jurídico de las Administraciones públicas y el régimen Estatutario de los funcionarios. Unificando el tratamiento de todos ellos.

        El TREBEP respeta la autonomía de las CCAA, que deben desarrollar sus propias leyes partiendo del mínimo común. En la Disposición Final Cuarta se establece los bloques relativos a la carrera administrativa, evaluación del desempeño y retribuciones complementarias. Se establecen los principios de actuación, operando como principios del ordenamiento: el servicio objetivo a los ciudadanos y al interés general, igualdad mérito y capacidad en el acceso y promoción profesiona y el sometimiento pleno a la ley y al derecho, igualdad de trato entre hombres y mujeres y la eficacia, transparencia y responsabilidad en la gestión de recursos humanos.

        El Ámbito de aplicación es universal para el sector público, extendiéndose a la AGE, CCAA, Ciudades Autónomas, Entidades Locales, Organismos públicos, Agencias dependientes y Universidades Públicas. Aunque cierto personal tiene un régimen de singularidades: el docente y el sanitario se rigen por su legislación específica, y por el TREBEP de forma supletoria.

        % Tabla del tema 024 p 19 y 20
        % \begin{table}[H]
        % \begin{tabularx}{\textwidth}{
        %       >{\raggedleft\arraybackslash}X
        %     | >{\raggedright\arraybackslash}X
        %     | >{\raggedright\arraybackslash}X
        %     }
        %     Título & Ámbito & Implicaciones para directivos TIC \\
        %     \midrule
        %     II    & Clases de Personal & Define la tipología de efectivos (Arts. 8 a 12: Funcionario de carrera, interino, laboral y eventual). Consagra la figura del Personal Directivo Profesional (Art. 13), clave para la gobernanza digital.\\
        %     III   & Derechos y deberes. Código de conducta & \\
        %     IV    & Adquisición y pérdida de la relación de servicio & \\
        %     V     & Ordenación de la actividad profesional & \\
        %     VI    & Situaciones administrativas & \\
        %     VII   & Régimen disciplinario & \\
        %     VIII  & Cooperación entre las Administraciones públicas & \\
        % \end{tabularx}
        % \end{table}

        \paragraph{La consolidacin Tecnolgica del modelo: Artículo 47bis y Derechos Digitales}
        La maduración digital de las administraciones públicas se formalizó mediante el Art. 47bis del TREBEP el cual transformó el teletrabajo de contingencia a modalizad organizativa. Este se asienta sobre los siguientes principios:
        \begin{description}
            \item[Definición y soporte tecnológico] La prestación de servicios a distancia fuera de las dependencias administrativas mediante el uso de las Tecnologías de la Información y la Comunicación. Se introduce el mandato imperativo deque la Administración debe proporcionar y mantener los medios necesarios para el desarrollo de la actividad, garantizando la seguridad de la información.
            \item[Voluntad y Reversibilidad] Tiene carácter estrictamente voluntario para el empleado y es reversible, tanto por instancia del trabajador, como por necesidades organizativas de la Administración, salvo situaciones excepcionales de fuerza mayor debidamente justificadas.
            \item[Gobernanza y Negociación colectiva] Los criterios objetivos de acceso a esta modalidad quedan supeditads a la negociación colectiva territorial. Debiendo garantizarse que las necesidades queden plenamente cubiertas. Su implantación tiene un impacto presupuestario neutro, no alterando las retribuciones del puesto.
            \item[Gestión por Objetivos] El estatuto vincula el teletrabajo a una reforma del modelo de gestión, determinado que la modalidad debe contribuir a una mejor organización laboral a través de la determinación de objetivos y la evaluación de su cumplimiento. Rompiendo con la cultura del presencialismo.
            \item[Blindaje de los derechos digitales] La adición del Art. 14.j bis (Por LO 3/2018 de Protección de datos) consagra el derecho a la intimidad frente al uso de dispositivos digitales, la monitorización por videovigilancia o geolocalización, y el derecho a la desconexión digital fuera del horario laboral.
        \end{description}

    \subsubsection{La regulación de la función públican en la AGE} % 19
        El TREBEP fija el denominador común, perocada AAPP debe articular su propia normativa de despliegue. Para la AGE, obligó a mantener un régimen de transición amparado en la Disposición Transitoria Final Cuarta, compuesto de dos normas fundamentales, la ley 30/1984 de Medidas para la REforma de la Función Pública con aspectos como la ordenación de los puestos de trabajo y niveles de funcionarios. Y el RD 364/1995 del Reglamento General de Ingreso del Personal al servicio de la AGE. Siendo la norma reglamentaria de referencia diaria para la gestión de concursos, libre designación y convocatorias de oposiciones estatales.

        En el caso de las TIC se complementa con directrices de la Secretaría de Estado de Función Pública e instrumentos de la Comisión de Estrategia TIC. Se define cómo las estructuras de personal deben adaptarse a marcos como el Esquema Nacional de Seguridad o el Esquema Nacional de Interoperabilidad.

        El RD-L 6/2023 supuso una gran reforma. Su título II constituye la Ley de Función pública, adaptando las bases del TREBEP al entorno digital con 4 grandes vectores de transformación.

        \paragraph{Planificación Estratégica de los Recursos Humanos basada en datos y competencias}
            Impuesta en un cambio metodológic: la gestión de personal pasa de trámite presupuestario a ser predictiva y orientada a datos. Hay planes Estratégicos de Recursos Humanos con planes plurianuales vinculando a plantillas con objetivos de digitalización y políticas públicas. Y Modelos de competencias, ya no solo por memorización, sino por mapas de competencias profesionales y habilidades blandas.

        \paragraph{Acceso y modelo de oposiciones mas ágiles y reducción de plazos de gestión}
            Se busca acelerar los procesos selectivos, reduciendo los tiempos de convocatoria y resolución para evitar perder candidatos cualificados. Profesionalización de tribunales, conórganos de selección permanentes formados por personal especializado y con criterios de paridad y objetividad.

        \paragraph{Evaluación del desempeño obligatoria ligada a retribución y carrera horizontal}
            Evaluación del desempeño obligatoria, periódica y universal para todo el personal de la AGE. Tiene efectos jurídicos directos, con impacto en la carrera horizontal, perceptción de retribuciones complementarias y en la continuidad en el puesto obtenido por concurso.

        \paragraph{Consolidación del Personal Directivo Profesional por Objetivos}
            Desarrolla el Art. 13, regulando el estatuto propio para puestos (como las Subdirecciones Generales TIC) que ejercen alta gerencia. Se designa de acuerdo a mérito, capacidad e idoneidad formal y funciones sujetas a evaluación y cumplimiento de objetivos.

\subsection{Clases de personal al servicio de las administraciones públicas y el personal directivo profesional} % 22
    % \subsubsection{Tipología de empleo público y clases de personal} % 22
    En las Admnistraciones públicas conviven diferentes tipos de empleados con regímenes jurídicos distintos. Esto garantiza el equilibrio y la flexibilidad de la Administración.

    Son empleados públicos quienes desempeñan funciones retribuidas en las Administraciones públicas al servicio de los intereses generales. De los cuales hay 4 categorías:

        \paragraph{Funcionarios de Carrera} 
            Es el núcleo duro y preferente del empleo público. se nombran legalmente tras la superación de un proceso selectivo y se incorporan de forma permanente. La relación no es contractual, sino estatutaria y regida por el derecho administrativo. Esto otorga inamovilidad en el cargo, destinado a garantizar su objetividad e imparcialidad en el cargo. Les corresponden a los funcionarios las potestades públicas y la salvaguardia de los intereses generales.

            Se organizan en Cuerpos y Escalas creadas por Ley y adscritos por grupos de clasificación según titulación exigida:
            \begin{description}
                \item[Grupo A] con los subgrupos A1 y A2. Exige un Grado Universitario. Cada Subgrupo corresponde al nivel de responsabilidad y dificultad de las pruebas.
                \item[Grupo B] Exige el título de Técnico Superior(Formación profesional de Grado Superior)
                \item[Grupo C] Con subgrupos C1 y C2, Exige Bachiller o Técnico para C1 y Educación Secundaria para C2
            \end{description}

        \paragraph{Funcionarios Interinos}
            Son nombrados por razones justificadas de necesidad y urgencia para desempeñar funciones propias de los funcionarios de carrera cuando se den circunstancias transitorias tasadas por la ley. Se rigen por el Derecho Administrativo y tienen los mismos derechos y deberes que los funcionarios de Carrera, salvo aquellas que interfieran con su naturaleza temporal.

            Los supuestos de nombramiento son únicamente los siguientes:
            \begin{description}
                \item[Existencia de plazas vacantes (Art. 10.1a)] Cuando no sea posible cubrirlas por funcionarios de carrera, con un plazo máximo de 3 años, tras los cuales debe ser cubierta por un Funcionario o ser amortizada.
                \item[Sustitución transitoria de tirulares (Art 10.1b)] Para cubrir bajas, licencias o reservas de plaza.
                \item[Ejecución de programas de caracter temporal (Art 10.1c)] Con un plazo máximo de 3 años ampliable un cuarto mas por las leyes de función pública territorial
                \item[Exceso o acumulación de tareas (Art 10.1d)] Por un plazo máximo de 9 meses dentro de un periodo de 18.
            \end{description}

        \paragraph{Personal laboral}
            Su vínculo con la administración es a través de un contrato de trabajo, quedan sujetos a la legislación laboral, con las modulaciones que impone el TREBEP, pueden ser fijos, indefinidos o temporales.

            Las leyes determinan qué puestos pueden ser desempeñados, respetando siempre la reserva de potestades al funcionario. Su campo suelen ser oficios, puestos instrumentales, asistenciales o gestión técnica no autoritaria. Su acceso a puestos fijos exgige procesos selectivos en convocatoria pública con principios de igualdad, mérito y capacidad.

            Los indefinidos no fijos son una creación jurisprudencial ante las irregularidades de la contratación temporal, otorga al trabajador el derecho a permanecer en el puesto hasta que sea provisto o amortizado.

        \paragraph{Personal eventual}
            Comunmente personal de confianza o asesoramiento político. Son empleados nombrados y cesados con total libertad de forma no permanente para funciones de confianza o asesoramiento especial. Tienen prohibido realizar funciones reservadas a los funcionarios de carrera.


    \subsubsection{Temporalidad y procesos de estabilización del empleo público} % 24
        Las figuras de la interinidad y contratación temporal han sido abusadas, tensionando el principio de mérito y provocando pronunciamientos del Tribunal de Justicia de la Unión Europea. Esto se respondió con la ley 20/2021 de medidas urgentes para la reducción de la temporalidad en el empleo público. Por la cual se establecen objetivos y mecanismos.
        \begin{description}
            \item[Tasa de temporalidad objetivo]Fijado en una temporalidad estructural inferior al 8\% en el conjunto de las administraciones públicas
            \item[Mecanismos extraordinarios de Estabilización] Autorización de convocatoria excepcional de plazas ocupadas temporal e ininterrumpidamente por dos vías, concurso oposición para plazas con continuidad de 3 a 5 años y el de concurso de méritos para plazas ocupadas ininterrumpidamente desde antes del 20260101.
            \item[Régimen de Sanciones y Disuasión] Para evitar futuras recaidas en precarización, con compensaciones exonómicas para el personal que cese por no superer el proceso, y la nulidad de pleno derecho de todo acto que suponga superar los plazos máximos de permanencia establecidos.
        \end{description}

    \subsubsection{El personal directivo profesional} % 24
        Con categoría singular del personal anterior, el Art.13 diseña el estatuto del Personal Directivo Profesional, configurándolo como una pieza innovadora del empleo público. Esto corresponde a una Administración compleja, tecnificada y digitalizada con un estrato gerencial profesionalizado entre las decisiones políticas del Gobierno y la ejecución técnica de las unidades administrativas. 

        El Estatuto recomienda al Gobierno del Estado y de las Comunidades un régimen jurídico específico que pivota sobre cuatro ejes:

        \begin{description}
            \item[Concepto Funcional] Este personal desarrolla funciones directivas profesionales.
            \item[Designación Competitiva] Atiende estrictamente a principios de mérito, capacidad e idoneidad. Garantizando un procedimiento público y concurrente entre candidatos.
            \item[Evaluación del Desempeño y Rendimiento] Sometido a un régimen de evaluación contínua de eficacia, eficiencia, responsabilidad de la gestión y control de resultados.
            \item[Exlcusión de la negociación colectiva] Siendo cúspide de la gobernanza y gestión, quedan excluidos de  la negociación colectiva con organizaciones sindicales.
        \end{description}

\subsection{Planificación de recursos humanos, Oferta de empleo público y acceso al empleo público} % 26
    La administración pública necesita gestionar a sus profesionales mediante herramientas de gestión. El legislador organiza el márco legal básico del \textbf{Título V} en tres capítulos:
    \begin{description}
        \item[Capítulo I] PLanificación de los recursos humanos.
        \item[Capítulo II] Estructura del empleo público.
        \item[Capítulo III] Provisión de puestos de trabajo y movilidad.
    \end{description}
    Siendo la eficacia de la planificación y gestión de los recursos uno de los principios de la actuación del empleo público.

    \subsubsection{La planificación estratégica de los recursos humanos} % 26
        Es el proceso que permite adaptar el personal disponible a las necesiades reales. Se pretende garantizar la eficacia y eficiencia de la prestación de servicios públicos y los recursos disponibles.

        Esto se realiza con
        \begin{description}
            \item[Plantillas de personal] qué plazas están dotadas
            \item[Relaciones de Puestos de Trabajo] Documentos púb licos que estructuran la organización de cada puesto, grupo, cuerpo, escala, provisión y retribuciones.
            \item[Planes de Ordenación de Recursos Humanos] Normas que permiten aplicar medidas como análisis de perfiles profesionales, modificación de estructuras, movilidad y formación. La aprobación de estos planes pide negociación con las organizaciones sindicales.
        \end{description}

        Los siguientes capítulos del Titulo V establecen (Capítulo II) la agrupación de los empleados en cuerpos, escalas y categorías y (Capítulo III) Los sistemas de carrera y provisión según méritos y excepcionalmente la libre disposición, con respeto a los derechos de movilidad voluntaria e intradministrativa. Esto ha pasado de ser una planificación burocrática a una estratégica.
        
        Se establece que la ordenación del empleo público debe hacerse de manera plurianual con planes de ámbito general, específicos sectoriales y de reestructuración. Diseñados a partir de mapas de competencias y perfiles profesionales. También exige un informe de ejecución cualitativo.

        El Art.71 obliga a crear Registros de Personal integrados en cada administración. Se determina que dichos registros deben tener en cuenta las particularidades de determinados colectivos (personal docente, fuerzas armadas o personal sanitario) y exige interoperatibilidad entre los registros de las Administraciones Públicas. Estatalmente se crea el Registro CEntral de Personal que insribe a los empleados públicos y anota los actos de su vida administrativa.

    \subsubsection{Instrumentos: La oferta de empleo público (OEP)} % 27
        Las necesidades de nuevo ingreso se operan mediante las Ofertas de Empleo Público. El régimen jurídico establece las siguientes reglas:
        \begin{description}
            \item[Aprobación y publicacioón] anualmente por los órganos correspondientes y publicada en el BOE
            \item[Obligación de convocar] permitiendo añadir un 10\% adicional para garantizar la cobertura de vacantes
            \item[Plazo de ejecución] improrrogable de 3 años.
        \end{description}

        La improrrogabilidad se hizo drástica con las medidas de limitación de la interinidad. La OPE se diseña entorno a tres pilares:
        \begin{itemize}
            \item Transformación digital
            \item Mejora de servicios esenciales
            \item Resiliencia tecnológica
        \end{itemize}

    \subsubsection{El Acceso al empleo público} % 28
        Para garantizar el derecho del ciudadano al acceso al empleo público, se obliga a las administraciones públicas y a sus órganos al sometimiento de los siguientes principios rectores:
        \begin{itemize}
            \item Publicidad
            \item Transparencia
            \item Objetividad e imparcialidad
            \item Independencia
            \item Adecuación
            \item Agilidad
        \end{itemize}

        Los aspirantes deben cumplir obligatoriamente los siguientes requisitos generales:
        \begin{description}
            \item[Nacionalidad Española] Con excepciones para los ciudadanos comunitarios y familiares.
            \item[Capacidad] para el desempeño de las tareas del puesto
            \item[Edad] Mínimo 16 años y menos de la edad máxima de jubilación forzosa
            \item[Habilitación] No haber sido separado mediante expediente disciplinario, ni inhabilitado por resolución judicial firme
            \item[Titulación] Según el grado al que se opte.
        \end{description}

        También se reflejan reglas de protección y equidad social. Regulando el acceso de funcionarios españoles de organismos internacionales y convalidando los conocimientos ya demostrados en puesto de origen. También regula una reserva de plazas no inferior al 7\% para personas con discapacidad, un 5\% para discapacidad física, psíquica y sensorial y un 2\% para discapacidad intelectual.

    \subsubsection{Sistemas selectivos, Órganos de selección y adquisición de la relación de servicios} %29
        Se regulan tres sistemas para la selección del personal funcionario. La \textbf{Oposición}, consistente en la superación de pruebas teóricas y prácticas objetivas. El \textbf{Concurso}, consistente en la valoración de méritos, tiene un caracter excepcional para supuestos singulares y tasados. Y el \textbf{Concurso Oposición}, donde se realiza una oposición, seguida de una fase de valoración de méritos.

        Los procedimientos podrán incluir la superación de un curso selectivo o un periodo de prácticas de caracter obligatorio y eliminatorios. En el grupo A1 estos cursos se imparten por centros de formación del INAP.

        La evaluación y calificación de las pruebas se encomienda a órganos de selección o tribunales. Deben ser imparciales y profesionalmente meritorios. Debe tener paridad de género y poseer titulación académica igual o superior a la exigencia. No pueden formar parte el personal de elección, interinos, eventual de confianza u organizaciones sindicales.

        Los funcionarios tienen cuatro requisitos para adquirir la condición de funcionario de carrera:
        \begin{description}
            \item[Superación del proceso selectivo]
            \item[Nombramiento formal]
            \item[Acto de acatamiento] expreso de la constitución y del estatuto de autonomía
            \item[Toma de posesión]
        \end{description}

        La relación es permanente, aunque se puede extinguir por las siguientes causas:
        \begin{itemize}
            \item Renuncia
            \item Pérdida de la naciuonalidad
            \item Jubilación
            \item Sanción disciplinaria
            \item Pena de inhabilitación absoluta o especial
        \end{itemize}

\subsection{Incompatibilidades del personal al servicio de las administraciones públicas} % 31
        El régimen garantiza, a través de las incompatibilidades, los principios de objetividad,imparcialidad y eficacia. Previniendo los conflictos de interés y asegurando la plena dedicación e independencia del personal público.

        El sistema prohibe las actividades que menoscaben los deberes del cargo, pero permite otras actividades de forma excepcional, tasada y autorizada. Por ello hay un código de conducta en dos dimensiones.
        
        El Art.53 detalla los Princípios Éticos con la abstención en caso de conflicto de intereses, neutralidad, confidencialidad, rechazo del trato de favor y de contraer obligaciones financieras o actividades que puedan colisionar con sus funciones públicas.

        El Art.54 Detalla Reglas que disciplinan el comportamiento diario, exigiendo el cumplimiento de la jornada y el horario, la ejecución de las tareas y el uso eficiente de los recursos públicos.


%     \subsubsection{Régimen general y fundamento ético} %31
%     \subsubsection{Marco normativo y ámbito de aplicación} % 32
%     \subsubsection{Incompatibilidad en el sector público y excepciones} % 32
%     \subsubsection{Incompatibilidad en el sector privado y actividades exentas} % 34

% 36